\documentclass[11pt,a4paper]{article}
\usepackage[utf8]{inputenc}
\usepackage[T1]{fontenc}
\usepackage[german]{babel}
\usepackage{amsmath,amssymb,amsthm}
\usepackage{mathrsfs}
\usepackage{geometry}
\geometry{a4paper, margin=2.5cm}

\newtheorem{theorem}{Satz}
\newtheorem{definition}{Definition}

\title{\textbf{Algebraische Geschlossenheit des $E_8$-Oktikkanals und die $\pi$-Leiter geradzahliger Gitterpackungen:\\Fano-Geometrie, $\text{Cl}(0,7)$-Spinorberechnung und das Nichtlineare Grad-10-Closure-Gate}}
\author{Thomas Hoffbauer \\ \small Bamberg, Deutschland}
\date{September 2026}

\begin{document}

\maketitle

\begin{abstract}
Wir untersuchen die algebraische Einbettung der klassischen Kugelpackungen in den geraden Dimensionen $2, 4, 6, 8$ im Zwischenraum $\mathbb{R}^7$ (imaginäre Oktonionen) mittels der Fano-Ebene und der normierten Divisionsalgebren. Es wird gezeigt, dass der betrachtete quadratische Fano-Metrikkanal keine Anisotropie trägt ($Q_{ij} = 0$). Innerhalb des $E_8$-Invariantensektors ist Grad 8 nach dem quadratischen Kanal der erste neue primitive Invariant. Über eine exakte ganzzahlige $\text{Cl}(0,7)$-Clifford-Darstellung und die Einbindung der primitiven $E_8$-Oktik-Invarianten nach Cederwall \& Palmkvist wird durch Cartan-Reduktion und den radialen $E_8$-Laplaceoperator ein vollständiger, nichtlinearer Abgleich bis zum Grad-10-Closure-Gate durchgeführt. Das System erweist sich als algebraisch in sich geschlossen, während eine photonische Ankopplung als offenes EFT-Gate markiert wird.
\end{abstract}

\section{Einleitung und der Kontext der $E_8$-Invarianz}
Der Ausnahmewert der Lie-Algebra $\mathfrak{e}_8$ und des zugehörigen $E_8$-Wurzelsystems beruht auf seiner einzigartigen Eigenschaft: Das $E_8$-Gitter ist das einzige gerade, selbstduale Gitter in acht Dimensionen. Während quadratische Invarianten (die Killing-Metrik) universell zugänglich sind, erfordert die Behandlung höherer Kasimir-Invarianten einen hochentwickelten Tensorformalismus. Wie Cederwall und Palmkvist gezeigt haben, liegen die primitiven Invarianten von $\mathfrak{e}_8$ bei den Graden $2, 8, 12, 14, 18, 20, 24$ und $30$. Der Oktik-Invariantentensor bildet dabei neben der Metrik den ersten neuen primitiven Kanal.

\section{Die $\pi$-Leiter geradzahliger Gitterpackungen}
Für die klassischen Gitterpackungen in den geraden Dimensionen $2, 4, 6, 8$ gelten
\[
\Delta_2=\frac{\pi}{2\sqrt{3}},\qquad
\Delta_4=\frac{\pi^2}{16},\qquad
\Delta_6=\frac{\pi^3}{48\sqrt{3}},\qquad
\Delta_8=\frac{\pi^4}{384}.
\]
Die aufeinanderfolgenden Potenzen von $\pi$ reflektieren zunächst den universellen Kugelvolumenfaktor
\[
V_{2m}(1)=\frac{\pi^m}{m!}.
\]

Darstellungstheoretisch besitzt $E_8$ zugleich eine maximale Rang-Untergruppe vom Typ $E_6\times SU(3)$, entsprechend $E_6\times A_2$ auf Wurzelsystemebene. Für die Adjungierte gilt
\[
\mathbf{248}=(\mathbf{78},\mathbf{1})\oplus(\mathbf{1},\mathbf{8})\oplus(\mathbf{27},\mathbf{3})
\oplus(\overline{\mathbf{27}},\overline{\mathbf{3}}).
\]
Die Rangzerlegung $6+2=8$ ist daher mit der additiven Struktur der $\pi$-Exponenten $3+1=4$ kompatibel.

Diese Parallelität darf jedoch nicht mit einer Produktzerlegung der Packungsdichten verwechselt werden. Tatsächlich gilt
\[
\Delta_2\Delta_6=\frac{\pi^4}{288},
\qquad
\Delta_8=\frac{\pi^4}{384}
=\frac{3}{4}\,\Delta_2\Delta_6.
\]
Die $\pi$-Leiter wird daher im Folgenden als geometrisch-darstellungstheoretische Korrespondenz, nicht als Ableitung der $E_8$-Packungsdichte aus $A_2$ und $E_6$, verstanden.

\section{Fano-Struktur und $\text{Cl}(0,7)$-Darstellung (Grad 2)}
Die Multiplikationsstruktur der 7 imaginären Oktonioneneinheiten im $\mathbb{R}^7$ wird durch den Strukturkonstanten-Tensor $c_{ijk}$ der Fano-Ebene beschrieben. Für den korrekt normalisierten Fano-Tensor gilt die Kontraktion:
$$C_{ij} := \sum_{m,p} c_{imp} c_{jmp} = 6\delta_{ij} \implies Q_{ij} = 0$$
Der betrachtete quadratische Fano-Metrikkanal trägt somit keine Anisotropie.

Um komplexe Zahlen zu vermeiden, wird die Clifford-Algebra $\text{Cl}(0,7)$ exakt über die oktonionische Links-Multiplikation realisiert: $\gamma_i = -L_{e_i}$. Mit einem $G_2$-invarianten Spinor $\eta$ reproduzieren die Spinorbilineare fehlerfrei die Fano-Drei- und Vierform:
$$\eta^T \gamma_i \gamma_j \gamma_k \eta = c_{ijk}, \quad \eta^T \gamma_i \gamma_j \gamma_k \gamma_l \eta = \psi_{ijkl} = (*\varphi)_{ijkl}$$

\section{Der Primitive $E_8$-Grad-8-Kanal und Cartan-Reduktion}
Über den Harish-Chandra-Homomorphismus schränkt sich der Cederwall-Palmkvist-Oktikterm $X_8^{\text{CP}}$ auf die Cartan-Unteralgebra ein:
$$X_8^{\text{CP}}\big|_{\mathfrak{h}, \, \beta=0} = \frac{1}{6} P_8^{\text{prim}}$$
wobei $P_8^{\text{prim}} = P_8 - \frac{1}{432000} P_2^4$ der Fischer-primitive Vertreter ist. Der auditierte, invariante Branching-Faktor des $E_8 \to F_4$-Restriktionskanals ist:
$$\sigma_8 = \frac{3}{4\sqrt{15}}, \quad \sigma_8^2 = \frac{3}{80}$$

\section{Das Nichtlineare Grad-10-Closure-Gate}
Die Kontraktion zweier Oktiktensoren kollabiert auf bekannte Strukturen:
$$t^{(A_1 \dots A_5 | BCD |} t^{|A_6 \dots A_{10})} {}_{BCD} = a \, t^{(A_1 \dots A_8} \delta^{A_9 A_{10})} + b \, \delta^{(A_1 A_2} \cdots \delta^{A_9 A_{10})}$$
mit den rationalen Koeffizienten $a = \frac{130065}{1778}$ und $b = \frac{1233765}{64516}$. Über den radialen $E_8$-Laplaceoperator auf der Cartan-Unteralgebra:
$$\mathcal{R}_{E_8} = \Delta_{\mathfrak{h}} + 2\sum_{\alpha>0} \frac{\alpha \cdot \nabla}{\alpha \cdot x}$$
wird das Polynom-Gate $K_{10} = \frac{\mathcal{R}_{E_8}^3 (F^2)}{8 \cdot (8 \cdot 7 \cdot 6)^2} \stackrel{!}{=} a F p_2 + b p_2^5$ durch Einsetzen der 120 positiven $E_8$-Wurzeln identisch erfüllt ($\text{residual} = 0$).

\section{Epistemische Trennung und offenes EFT-Matching}
Die vorangegangenen Resultate bestimmen einen kanonischen algebraischen Grad-8-Kanal. Sie bestimmen jedoch weder seine Kopplung an das elektromagnetische Feld noch eine dimensionsbehaftete Längenskala.

Ein schematischer, dimensionskonsistenter EFT-Ansatz kann in der Form
\[
\mathscr{K}_\gamma
=
\mathscr{K}_{\gamma,0}
+
\eta_8\,\ell_8^8\,
\mathscr{F}_8
+
\mathcal{O}\left((\ell_8 p)^{12}\right)
\]
geschrieben werden, wobei $\mathscr{F}_8$ ein homogenes Polynom vom Grad acht in den relevanten Frequenz- und Impulsvariablen bezeichnet.

Die Skala $\ell_8$ bleibt in der vorliegenden Arbeit unbestimmt. Entsprechend ist auch $E_8 = \frac{\hbar c}{\ell_8}$ nur eine Parametrisierung einer möglichen Energieskala und keine abgeleitete neue Naturkonstante.

Insbesondere wird hier weder eine Identifikation $\ell_8 \propto \ell_P$ vorausgesetzt noch ein bestimmter numerischer Faktor festgelegt. Ebenso bleiben die gaugeinvariante Abbildung $t_8 \longrightarrow \mathcal{O}^\gamma_8$ und eine daraus folgende Casimir-Korrektur Gegenstand weiterer Untersuchungen.

\end{document}